Подытожу.

Документация пишется на английском.
Документируются только объекты которые имеют сигнатуры.

Формат ссылок:

Simple Grammar:

  | - alternation

link:
  | "{@link" /\s+/ link-ref "}"

link-ref:
  | entity
  | package-path entity

package-path:
  | identifier "/" path "#"
  | identifier "#"

path:
  | identifier "/" path
  | identifier

entity:
  | identifier "." entity
  | ref

ref:
  | identifier
  | "(" identifier ":" ref-selector ")"

ref-selector:
  | ref-type
  | ref-index

ref-type:
  | "class"
  | "enum"
  | "function"
  | "interface"
  | "namespace"
  | "type"
  | "variable"

ref-index:
  | /[1-9][0-9]?/

identifier:
  | /^[_$\p{IDS}][$\p{IDC}]*$/u





@link <package> <path> # <name>

Unable to resolve reference "(UnitSet:constructor)": Unsupported system selector "constructor"


/**
 * Это ссылка {}
 */

// For merged declarations:
// 'class',
// 'enum',
// 'function',
// 'interface',
// 'namespace',
// 'type',
// 'variable'

<type>:
 type
 enum
 constant
 variable



type-param:

/**
 * ...
 *
 * @typeParam <name> - Description. 1-2 sentences.
 */

Если в тексте встречается упоминание параметра, оно оборачивается в `...`.

/**
 * ... constructs a new type from the passed `<name>` by ...
 */

type:

/**
 * The {@link <name>} type represents / constructs a new type / do something / ...
 *
 * @remarks
 *
 *
 * @typeParam ...
 *
 * @typeParam ...
 */


Если в тексте встречается упоминание параметра, оно оборачивается в `...`.

интерфейсы
+ сигнатуры generic параметров

@typeParam Name - Description

классы
сигнатуры generic параметров

@typeParam Name - Description

модули, константы (или переменные) верхнего уровня (в рамках файла или модуля),  сами модули (пока про них ничего не написано, будет потом).
свойства
